% !Mode:: "TeX:UTF-8"

\chapter{多模型协同进行的可信问答}

\section{引言}

二三两章分别从错误识别和错误缓解两个角度讨论了大语言模型在复杂问答任务中的可靠性问题。
错误识别强调对模型生成过程进行诊断，错误缓解则尝试通过反思验证数据和强化学习改进模型本身。
仅依赖单一生成模型进行问答系统的构建仍然存在验证方面的不足。

本章面向多模型协同场景，构建一种带验证的可迭代可信问答流程。
该流程以策略模型负责候选答案生成，以验证模型负责质量检查和反馈生成，
并通过标准答案匹配对候选结果进行基础筛选。对于未通过质量检查的样本，
系统将验证反馈重新输入生成模型，引导模型在下一轮中进行反思式自搜索。
与一次性生成答案相比，该方法能够持续利用过去回复和验证反馈，
在多轮迭代中逐步增加可用样本数量，从而提高问答数据构建和答案生成的可信度。

本章实验基于 PopQA、MuSiQue、HotpotQA、Bamboogle 和 NQ 五个问答数据源展开。
实验结果表明，在最多 8 轮迭代、每轮 6 个候选答案以及验证模型 8 次判断的设置下，
系统最终在五个数据源上表现提升明显。
这说明多模型协同验证不仅能够判断答案是否命中标准答案，
还能够从推理、检索和回答一致性角度筛选出更适合用于后续训练或部署的可信样本。


\section{方法设计}

\subsection{总体框架}

本章方法的总体流程如图~\ref{fig:multi-model-workflow} 所示。
系统首先接收原始问题，由策略模型生成多个候选自检索答案。
每条候选轨迹包含思考、搜索、信息生成和最终回答四类内容。
随后，系统使用标准答案进行基础筛选，保留至少能够命中标准答案的候选。
在此基础上，验证模型对候选轨迹进行质量检查，判断其中是否存在推理、知识调用或答案整合错误。
若候选通过验证，则被保留为可用样本；若未通过验证且尚未达到最大迭代轮数，
则根据验证反馈重新生成候选答案。

\begin{figure}[htbp]
\centering
\includegraphics[width=0.96\textwidth]{multi_model_workflow}
\caption{多模型协同可信问答流程}
\label{fig:multi-model-workflow}
\end{figure}

该框架包含三个核心角色。
第一，策略模型负责生成候选回答，是系统探索答案空间的主体。
第二，标准答案匹配模块负责快速判断候选答案是否至少出现过完全正确结果。
第三，验证模型负责从过程层面评估候选轨迹，并为失败样本生成反馈。
通过三者协同，系统既保留了生成模型的探索能力，也引入了独立验证模型的质量约束。

\subsection{候选生成与基础筛选}

候选生成阶段采用多样化采样策略。
本章实验中，每轮为每个问题生成 6 个候选答案，生成温度设置为 0.7，
\texttt{top\_p} 设置为 0.8，\texttt{top\_k} 设置为 20。
该设置使模型能够在保持一定答案质量的同时探索多种查询和回答路径。
对于每个问题，只要候选集合中至少存在一个答案与标准答案完全匹配，
系统就将该问题记为在当前轮次达到累计完全匹配。

完全匹配指标只评价候选答案本身是否命中标准答案，
并不等价于完整生成过程是否可信。在最终答案上偶然命中标准答案的候选项可能在中间查询、生成信息或解释过程存在明显错误。
因此，完全匹配指标只能用于从候选集合中找出可能正确的样本，
这些样本还需要经过验证模型进一步检查。

\subsection{验证与反馈再生成}

验证阶段采用第二章确定的验证模型作为质量评估模型。
验证模型接收原始问题、模型完整输出和标准答案，
并按照“审查生成过程”和“解释错误原因”的步骤输出结构化判断。
若样本存在错误，验证模型需要给出错误描述以及反思反馈；
若样本正确，则说明各步骤为何合理。

本文提出一种用于衡量模型生成质量的指标\texttt{reject@k}。
\texttt{reject@k}指验证模型一次关于同一样本给出 k 个独立验证，
其中至少 1 个判断样本存在错误的概率。

\begin{align}
\mathrm{reject@}k
= 1 - \frac{\binom{\mathrm{c}}{k}}{\binom{m}{k}}
\end{align}

式中，\texttt{c}验证模型为判断样本正确的验证数目，\texttt{m}为验证模型采样次数。


\section{实验设置}

\subsection{数据集与运行配置}

实验覆盖 Bamboogle、HotpotQA、MuSiQue、NQ 和 PopQA 五个数据源。
其中 Bamboogle 包含 125 个问题，其余四个数据源各包含 500 个问题，
原始问题总数为 2125。

超参数设置见附录表~\ref{tab:appendix-multi-model-config}。

\subsection{评价指标}

本文使用两类指标评价迭代式可信问答流程。
第一类为累计完全匹配指标。
该指标表示从第 0 轮到当前轮为止，某个问题是否曾经生成过至少一个与标准答案完全匹配的候选答案。
第二类为累计验证通过指标，即从第 0 轮到当前轮为止，
某个问题是否曾经生成过通过验证模型质量检查的候选轨迹。

此外，本文报告验证通过率和验证/EM 比例。
验证通过率以数据源问题总数为分母，反映最终可用样本覆盖程度；
验证/EM 比例以累计完全匹配样本数为分母，
反映在已经命中标准答案的样本中，有多少比例能够进一步通过过程质量检查。
该比例越高，说明模型不仅更容易答对问题，而且更容易生成过程一致、可解释且可复用的回答。

\section{实验结果与分析}

\subsection{最终累计结果}

第 7 轮结束后的逐数据源准确率与验证率如表~\ref{tab:multi-model-rates-main} 所示，
完整计数和验证/EM 比例见附录表~\ref{tab:appendix-multi-model-final}。

\begin{table}[htbp]
\caption{不同数据源上的准确率与验证率}
\label{tab:multi-model-rates-main}
\vspace{0.5em}\centering\wuhao
\begin{tabular}{lcc}
\toprule
数据源 & 准确率 & 验证率 \\
\midrule
Bamboogle & 96.80\% & 80.80\% \\
HotpotQA & 89.40\% & 75.60\% \\
MuSiQue & 92.40\% & 69.60\% \\
NQ & 91.00\% & 67.40\% \\
PopQA & 98.60\% & 66.80\% \\
\bottomrule
\end{tabular}
\end{table}

每轮迭代后的累计结果见附录表~\ref{tab:appendix-multi-model-iteration}

从累计完全匹配结果看，五个数据源均取得较高覆盖率。
这说明多候选采样和反馈再生成能够显著提高模型至少一次答对问题的概率。
与完全匹配结果相比，验证通过结果更能体现可信问答的质量约束。
事实型问题上最终答案命中并不必然意味着生成过程可靠，
验证模型仍会过滤掉一部分过程存在错误或解释不充分的候选。


\subsection{结果讨论}

综合最终累计结果可以看出，多候选生成与反馈再生成能够显著提高标准答案命中率，验证反馈能够有效引导策略模型修正先前回复中的错误。
验证模型提供了比完全匹配更严格的质量控制，验证通过样本少于答案完全匹配样本。本章所述多模型协同进行的可信回复方法
可以应用于高质量回复的生成，如问答系统与高质量数据生成。

\section{本章小结}

本章研究了多模型协同进行的可信问答方法。
针对单一生成模型在自检索问答中难以同时保证答案正确性和过程可信性的问题，
本文构建了由策略模型、标准答案匹配模块和验证模型共同组成的迭代式问答流程。
策略模型负责生成多个候选自检索轨迹，标准答案匹配模块负责基础筛选，
验证模型则进一步检查推理、搜索、信息生成和最终回答是否存在错误。
对于未通过验证的样本，系统将验证反馈用于下一轮反思式再生成，
从而形成“生成-验证-反馈-再生成”的闭环。

实验结果说明，多模型协同验证能够在答案命中之外提供更严格的过程质量控制，
并为后续可信问答数据构建和模型训练提供可复用的技术路径。

% Local Variables:
% TeX-master: "../thesis"
% TeX-engine: xetex
% End:
